\begin{corollary}[$K_1$ in the translation-covariant scheme]
\label{cor:c11-K}
Under the hypotheses of Theorems \ref{thm:c8-adjoint} and \ref{thm:c11-arity2}, on a family of Definition \ref{def:c11-families} at the origin of one length, in the translation-covariant scheme at the values $\theta=\theta_\lambda=\theta_V=-\tfrac13$ of Theorem \ref{thm:c9-theta2}, $K_1$ along $K_{\bar b}$ equals $\hbar\big[\{S^{(1)}(K_{\bar b}),-\},\mathrm{Gram}_{\rm tree}^{-1}\mathrm{Gram}_{\rm 1}\big]$, the statement $K=P^\dagger$ of Corollary \ref{cor:c8-gram}. For equal coefficients $\theta_P$ along $P_b$, $K_1$ along $K_{\bar b}$ is $(\theta-\theta_P)$ units times $\{S^{(1)}(K_{\bar b}),-\}$ plus the same commutator, with the radial Gram matrix of Definition \ref{def:c8-inner} in that scheme. On a scalar bilinear at the origin both sides vanish: the defect cycles give $K_1(O)=0$ there (Proposition \ref{prop:c11-SKbil} (ii)), and the commutator vanishes on a primary at the origin on which $\mathrm{Gram}_{\rm tree}^{-1}\mathrm{Gram}_{\rm 1}$ is diagonal. On its descendants $\partial_\mu O$ the relation equates the intrinsic kernels $\Upsilon^0_{K_b}$ of the descendant pairs with the difference of the two-point normalisations in Weyl weights $w_O+1$ and $w_O$ times the classical transformation. The component $K_1$ itself is computed from its own defect cycles at a finite endpoint and from the action-side solution along $K_b$ (Proposition \ref{prop:c11-SKbil} (ii)); this corollary is a cross-check of it through the inner product and enters no number.
\end{corollary}
